\documentclass[11pt]{article}
\usepackage{mathunal-base}
\mathunalfuente{serif}
\newcommand{\rp}{\underline{\hspace{3cm}}}

\renewcommand{\examtitulo}{Ecuaciones Diferenciales (1000007) --- Primer examen parcial supletorio}
\renewcommand{\examfecha}{Semestre 01-2015, 16 de Marzo de 2015}

\begin{document}

\examheader

\vspace{4pt}
\noindent\textit{Duración: 1 hora y 40 minutos. Instrucciones: No está permitido el uso de calculadora ni de cualquier otro tipo de aparato electrónico, incluyendo teléfonos móviles, los cuales deben permanecer apagados durante el tiempo de duración del examen. Tampoco está permitido hacer preguntas a las personas encargadas de la vigilancia.}

\vspace{6pt}
\noindent\textbf{1. (25\%)} Resuelva la ecuación diferencial
\[
y\,dx+\left(2xy-e^{-2y}\right)dy=0.
\]

\vspace{6cm}

\newpage

\noindent\textbf{2. (25\%)} Obtenga la familia de curvas ortogonales a la familia
\[
y^3+3x^2y=C.
\]

\vspace{7cm}

\newpage

\noindent\textbf{3.} Un tanque contiene inicialmente 500 galones de agua con 25 libras de sal en solución. Se hace entrar agua limpia al tanque a razón de 2 galones por minuto y se deja que la mezcla salga del tanque a razón de 2 galones por minuto. La mezcla que sale de este primer tanque se vierte a un segundo tanque que tiene una capacidad de 100 galones y que inicialmente contiene 50 galones de agua pura.

\vspace{0.2cm}
\noindent\textbf{a) (15\%)} Encuentre la cantidad de sal en el primer tanque en cualquier instante.

\vspace{3.5cm}

\noindent\textbf{b) (15\%)} Encuentre la cantidad de sal en el segundo tanque en el instante que la solución empieza a derramarse.

\vspace{3.5cm}

\newpage

\noindent\textbf{4. (20\%)} Halle la solución explícita del problema de valor inicial
\[
x^2\frac{dy}{dx}=y-xy, \qquad y(-1)=-1,
\]
y determine el intervalo abierto más grande en el cual dicha solución es válida.

\vspace{7cm}

\vfill
\rule{\linewidth}{0.4pt}\\[1pt]
{\scriptsize\textit{Transcripción digital --- MathUNAL}\hfill\textcolor{unalblue}{\textbf{1}}}

\end{document}
